
\section{Introduction}

Knowledge evolution is a central question in information science. Existing science mapping research has demonstrated how knowledge structures emerge, transform, and diffuse through scholarly communication networks. However, technological knowledge evolution remains insufficiently understood.

Patents are often treated as innovation evaluation indicators, such as patent counts, citation impact, or technology output measures. Such approaches provide useful measurements of innovation performance but overlook the evolutionary processes embedded in technological knowledge systems.

This study argues that patents should be conceptualized as structured technological knowledge objects. They record not only technological outputs but also the processes through which knowledge is acquired, recombined, transformed, and diffused.

Accordingly, this study develops a technological knowledge evolution framework that models innovation as a dynamic process:

\begin{align*}
\text{Knowledge Acquisition} &\rightarrow \text{Knowledge Recombination}\\
&\rightarrow \text{Knowledge Emergence}\\
&\rightarrow \text{Knowledge Diffusion}\\
&\rightarrow \text{Knowledge Transformation}
\end{align*}

Using large-scale patent data from Chinese universities, this study investigates how technological knowledge evolves and how universities follow different evolutionary trajectories.

The contributions are threefold.

First, this study develops a knowledge evolution framework that shifts innovation analysis from static patent evaluation toward dynamic evolutionary mechanisms.

Second, this study extends science mapping from scientific knowledge systems to technological knowledge systems by integrating technological, semantic, and citation relationships.

Third, this study reveals heterogeneous technological evolution trajectories and provides a basis for differentiated innovation governance.
